\SourceLeaf{288}%
\begin{center}
\begin{tabular}{lr}
Brocken&\(+0''00141\)\\
Hohehagen&\(-0{,}00094\)\\
Inselsberg&\(-0{,}00511\)
\end{tabular}
\end{center}
\hspace*{\parindent}Ein zweites Beispiel entlehne ich aus der trigonometrischen Vermessung der Schweiz\footnote{Ergebnisse der trigonometrischen Vermessungen in der Schweiz, herausgegeben von J. \textsc{Eschmann}. Zurich 1840. S. 79. 99. 189. 190. 196.}, wo das grösste Hauptdreieck zwischen den Punkten Chasseral, Suchet, Berra eben an die Grenze der Ausdehnung unserer Hülfstafel fällt. Wir haben für diese Punkte
\begin{center}
\begin{tabular}{l|c|c|r}
&\multicolumn{2}{c|}{Breite.}&\\
&auf dem Ellipsoid&auf der Kugel&\(k\)\\\hline
Chasseral&\(47^0\quad8'\quad1''\)&\(47^0\quad6'\quad33''\)&\(6''137\)\\
Suchet&\(46\quad46\quad23\)&\(46\quad44\quad57\)&\(6{,}948\)\\
Berra&\(46\quad40\quad36\)&\(46\quad39\quad11\)&\(7{,}173\)
\end{tabular}
\end{center}
Die Logarithmen der Dreiecksseiten in Metern sind
\begin{center}
\begin{tabular}{lr}
Suchet-Berra&\(4{,}7474503\)\\
Berra-Chasseral&\(4{,}7133766\)\\
Chasseral-Suchet&\(4{,}7808768\)
\end{tabular}
\end{center}
Die Azimuthe
\begin{center}
\begin{tabular}{lrrr}
&\multicolumn{3}{c}{Standpunkt Chasseral}\\
Suchet&\(48^0\)&\(36'\)&\(41''\)\\
Berra&349&21&54\\
&\multicolumn{3}{c}{Standpunkt Suchet}\\
Chasseral&228&10&40\\
Berra&280&47&19\\
&\multicolumn{3}{c}{Standpunkt Berra}\\
Suchet&101&18&40\\
Chasseral&169&27&22
\end{tabular}
\end{center}
\hspace*{\parindent}Hieraus ergeben sich die Reductionen der Sphaeroid-Azimuthe auf die Kugel-Azimuthe
